\title{Byte Latent Transformer: Patches Scale Better Than Tokens}

\begin{abstract}
This work presents the Byte Latent Transformer ({\textsc BLT}), a novel byte-level large language model architecture which, for the first time, achieves performance at par with tokenization-based LLMs at scale, while also delivering substantial gains in inference speed and resilience. {\textsc BLT} represents input bytes as adaptively sized patches that become the core processing units. The segmentation of these patches leverages the entropy of the subsequent byte, dedicating greater computational and model resources to regions of higher data complexity. We conduct the inaugural scaling study of byte-level models with flop control, scaling up to 8B parameters and training on 4T bytes. Our findings establish that models can effectively scale directly on raw byte sequences, removing the dependency on a static vocabulary. Both model training and inference see increased efficiency through the adaptive selection of longer patches in more predictable contexts. Qualitatively, the approach enhances reasoning capability and generalization over rare cases. In summary, when inference budget is fixed, {\textsc BLT} surpasses tokenization-dependent models in scaling, driven by the simultaneous expansion of patch size and overall model capacity.
\end{abstract}

\section{Introduction}
\label{section:intro}

We present the Byte Latent Transformer~(\textbf{BLT}), an innovative architecture that dispenses with tokenization by directly processing raw byte sequences. For the first time, BLT achieves parity with traditional tokenization-dependent models on large-scale benchmarks, while also delivering marked gains in both computational efficiency and resilience (\S\ref{sec:robustness}).
Traditional large language models (llms) are generally trained in an almost entirely end-to-end manner, with the notable exception of the tokenization stage—a heuristic and static process that segments byte streams into discrete tokens. This preliminary step inherently influences text compression strategies and introduces several limitations, including sensitivity to linguistic domain or data modality~\citep{dagangetting}, increased vulnerability to noisy input~(\S\ref{sec:robustness}), an absence of explicit orthographic structure~\citep{edman2024cute}, and unfair treatment of less-represented languages~\citep{liang2023xlm,petrov2024language,limisiewicz2024myte}.

Tokenization has historically played a crucial role because training large language models on raw byte sequences directly has been prohibitively expensive at scale, largely due to the increased sequence lengths~\citep{xue2022byt5}. Earlier research has addressed this issue by adopting self-attention mechanisms with improved efficiency~\citep{el2020characterbert, clark2022canine}, or by exploring architectures that eschew attention entirely~\citep{wang2024mambabyte}~(\S\ref{section:related_work}). Nonetheless, such approaches predominantly offer benefits for training \textit{small models}. When scaling up, the principal computational burden in Transformer architectures is attributed to the extensive feed-forward network layers—applied to every byte—rather than to the attention mechanism itself.

We introduce a dynamic and adaptable approach for organizing bytes into \textit{patches}~(\S\ref{section:patching}) and present a novel model architecture capable of jointly processing information at both the byte and patch levels. In contrast to tokenization, {\textsc BLT} does not rely on a predefined patch vocabulary. Instead, arbitrary byte groupings are transformed into latent representations of patches through compact, trainable encoder and decoder components. Our results demonstrate that this framework enables a \textit{greater} efficiency in compute allocation compared to models grounded in tokenization.

Large language models based on tokenization distribute an equal amount of computation to each token, regardless of its inherent complexity. This approach prioritizes consistent performance but compromises on computational efficiency, as tokenization relies on compression methods that do not always align with prediction difficulty. Our architecture is grounded in the principle that computation should be adaptively distributed according to demand. For instance, using a highly complex transformer is often unnecessary when predicting the conclusion of most words, as these are typically straightforward, low-uncertainty choices, especially in contrast to selecting the initial word of a sentence, which tends to be more challenging. In the design of {\textsc BLT}~(\S\ref{section:architecture}), this philosophy manifests as three transformer modules: two compact local models operating at the byte level, and one substantially larger global latent transformer~(\autoref{fig:arch_high_level}). 

{\textsc BLT} achieves dynamic compute distribution by segmenting sequences into patches. This partitioning is informed by the entropy of next-byte predictions, ensuring that each patch contains bytes characterized by similar levels of information content.

We conduct the first study on flop-controlled scaling for byte-level architectures, exploring models as large as 8B parameters and trained on up to 4T bytes, and demonstrate that large-scale, end-to-end training from raw bytes is feasible without relying on fixed-vocabulary tokenization. In terms of training, {\textsc BLT} achieves performance at parity with Llama 3 under flop-equivalent settings,\footnote{We quantify the computational requirements by summing the Floating Point OPerations (flops) a model consumes.} but is able to reduce inference flop usage by as much as 50\% (\S\ref{section:scaling}). 

Additionally, we provide evidence that direct byte-level modeling yields clear advantages for capturing rare patterns within the data, with {\textsc BLT} exhibiting increased resilience to noise in comparison to traditional tokenizer-dependent models, and demonstrating stronger character-level competencies as shown in tasks relating to orthographic and phonological processing, as well as low-resource machine translation (\S\ref{sec:robustness}). 

Moreover, {\textsc BLT} enables scaling of both model size and patch size concurrently while keeping the inference flop cost unchanged. Increasing the patch size results in computational savings since the global latent transformer runs less frequently, allowing resources to be redirected toward model scaling. Through scaling studies at a fixed inference-flop budget (\autoref{fig:fixed_inference_scaling}), we find that such models exhibit markedly superior scaling behavior relative to models built around tokenization.

To conclude, the primary contributions of this work are as follows: 1) We present {\textsc BLT}, a byte latent llm framework that adaptively distributes computational resources to boost flop efficiency; 2) We demonstrate that our approach matches the training flop utilization of Llama 3 up to the 8B parameter scale, with the capability to exchange small degradations in evaluation benchmarks for as much as a 50\% gain in flop efficiency; 3) {\textsc BLT} models introduce a novel pathway for llm scaling, allowing for increased model capacity under a constant inference cost; 4) Our results indicate that {\textsc BLT} exhibits enhanced resistance to input perturbations and better captures sub-word characteristics in the input, capabilities that token-based llms often lack. The codebase for both training and inference with {\textsc BLT} has been made publicly available at~\url{https://github.com/facebookresearch/blt}.

\section{Patching: From Individual Bytes to Groups of Bytes}
\label{section:patching}

Dividing bytes into discrete units called \textit{patches} enables {\textsc BLT} to flexibly allocate computational resources according to contextual needs. As illustrated in Figure~\ref{fig:patching_types}, there are multiple approaches for grouping bytes into patches. More formally, a patching function $f_p$ transforms a byte sequence $\pmb{x} = \{x_i,|i=1,\ldots n\}$ of length $n$ into a series of $m < n$ patches $\pmb{p} = \{p_j|j=1,\ldots,m\}$, assigning each $x_i$ to either 0 or 1, with 1 marking the beginning of a patch. Regardless of whether a model is token- or patch-centric, the main driver of computational load remains the number of sequential operations performed by the central Transformer block. Within {\textsc BLT}, this computational effort corresponds to the quantity of patches into which the data is partitioned by a particular patching function. Therefore, the typical length of a patch, referred to as the \textit{patch size}, dictates the resources required to process input data, both during training and at inference time, for any selected patching strategy~(\S\ref{section:flops}).

We proceed by outlining three kinds of patching functions: segmentation using a constant number of bytes per patch~(\S\ref{section:static-patch}), whitespace-based patching~(\S\ref{section:white-patch}), and patching adaptively via entropy estimates from a compact byte language model~(\S\ref{section:dyn-patch}). Additionally, we examine methods for incremental patching and elucidate distinctions between patching and conventional tokenization~(\S\ref{section:bpe-patch}).

\subsection{Strided Patching Every K Bytes}
\label{section:static-patch}

One of the simplest methods for aggregating bytes is to partition them into uniform patches of size $k$, as demonstrated in MegaByte~\citep{yu2023megabyte}. Employing a fixed stride simplifies both implementation during training and deployment at inference, offers a transparent approach for adjusting average patch size, and thereby facilitates control over computational requirements. Despite these advantages, this patching strategy is associated with notable drawbacks. Primarily, the allocation of computational resources does not adapt dynamically to the areas where it is most required: a given transformer iteration $j$ may be wasted when dealing with less informative content such as whitespace in code, or might under-serve segments rich in informational content, like mathematical expressions. Additionally, this approach introduces variability and a lack of contextual coherence in how similar byte sequences are segmented; for instance, identical words may end up fragmented inconsistently across different patches.

\subsection{Space Patching}
\label{section:white-patch}

\citet{slagle2024spacebyte} introduces an enhancement over the strided patching approach by generating new patches whenever a space-like byte\footnote{
    Space-like bytes encompass any byte that is neither a Latin letter, a digit, nor a UTF-8 continuation byte. Additionally, patches are required to include at least one byte that is not space-like. } is encountered, as these often represent intuitive linguistic boundaries across numerous languages. Under the Space patching strategy, every word is allocated a distinct latent transformer step (i.e., increased computational cost) for modeling, promoting consistent patching of words throughout different sequences and focusing computational resources on challenging prediction tasks, which frequently occur immediately after a space. As an illustration, predicting the initial byte in the answer to “Who composed the Magic Flute?\uline{\hspace{.6em}}” poses a greater challenge than forecasting subsequent bytes after the letter “M”, since identifying the first character dramatically narrows down possible continuations, whereas predicting the remainder of "Mozart" becomes relatively straightforward. Nevertheless, space patching has limitations: it does not accommodate every language and domain effectively and, more crucially, lacks adaptability in patch size. To address these issues, we now present a novel patching approach, leveraging the observation that the initial bytes of words usually present the greatest prediction challenge, and that inherently supports flexible control over patch size.

\subsection{Entropy Patching: Using Next-Byte Entropies from a Small Byte LM}
\label{section:dyn-patch}

Instead of employing a rule-based method like identifying whitespace, we adopt a data-driven strategy to detect points of elevated uncertainty in next-byte prediction. We propose \textit{entropy patching}, a technique that determines patch boundaries by leveraging entropy calculations.

We train a small byte-level auto-regressive language model on the training data for {\textsc BLT} and compute next byte entropies under the LM distribution $p_{e}$ over the byte vocabulary $\mathcal{V}$:

\begin{equation}
H(x_i) = \sum_{v \in \mathcal{V}} p_{e}(x_i=v|\pmb{x}_{<i}) \log p_{e}(x_i=v|\pmb{x}_{<i})
\end{equation}

We explore two approaches for determining patch boundaries based on entropies $H(x_i)$. The first approach selects points where the entropy exceeds a fixed global threshold, as shown in~\autoref{fig:patching}. The second method selects points whose entropy is significantly higher than that of the preceding point. This latter technique can alternatively be viewed as detecting locations where there is a departure from the generally monotonic decrease in entropy within a patch.

\begin{align*}
    \text{Global Constraint}\hspace{1cm}H(x_t)& > \theta_g\\
    \text{Approx. Monotonic Constraint}\hspace{1cm}H(x_t)& - H(x_{t-1}) > \theta_r
\end{align*}

Patch boundaries are determined through a simple preprocessing procedure carried out at the dataloading stage. This approach contrasts with that of~\citet{nawrot-etal-2023-efficient}, in which a classifier is employed to infer entropy-based patch boundaries. In our experiments~(\S\ref{section:experiments}), we evaluate and compare both strategies for separating low- and high-entropy bytes.

\subsection{The Byte-Pair Encoding (BPE) Tokenizer and Incremental Patching}
\label{section:incremental-patching}
\label{section:bpe-patch}

A large number of contemporary llms, such as our baseline Llama 3, rely on subword tokenizers like bpe~\citep{gage1994new, sennrich-etal-2016-neural}. Throughout this work, ``tokens'' denote groups of bytes selected from a \textit{finite} vocabulary established before training begins, whereas ``patches'' pertain to sequences that are grouped on the fly and are not limited to a predetermined set. A key distinction between these two representations is that, when using tokens, the model does not possess explicit access to the raw byte-level information.

An important advancement introduced by {\textsc BLT} in comparison to tokenization-based models lies in its revised balance between vocabulary size and computational requirements. Traditional large language models manage an inherent compromise: enlarging the vocabulary results in longer average tokens, which decreases the number of inference steps, yet simultaneously increases the output dimensionality of the model's final projection layer. Consequently, tokenization-based strategies offer limited potential for substantial shifts in both token length and computational overhead during inference. For instance, Llama 3 achieves an increase in average token size from 3.7 to 4.4 bytes, but this comes at the expense of enlarging its embedding table by a factor of four relative to Llama 2.

During generation, {\textsc BLT} must determine at each point in the byte sequence whether it has reached a patch boundary, as this influences whether additional computation is allocated to the Latent Transformer. This determination must be made without knowledge of subsequent, as-yet-unproduced bytes in the sequence. Consequently, the procedure for defining patches cannot rely on future bytes to segment the sequence. More precisely, a patching method $f_p$ is considered to be incrementally patching if it fulfills the following criterion:
$$f_p(\pmb{x}_{<i}) = f_p(\pmb{x})_{<i}$$
Notably, bpe does not qualify as an incremental patching method since identical prefixes may be tokenized in different ways depending on what follows, thus failing to satisfy the above property.\footnote{It is possible to modify bpe to act as an incremental patching scheme by introducing a dedicated boundary marker, but this approach lengthens the resulting byte sequence.}

\section{{\textsc BLT} Architecture}
\label{section:architecture}
The {\textsc BLT} framework consists of a primary global autoregressive language model that processes representations at the patch level, supported by two lightweight local components. These subsidiary models are responsible for transforming byte sequences into patch representations and for reconstructing bytes from patches, as depicted in~\autoref{fig:arch_high_level}.

\subsection{Latent Global Transformer Model}

\textit{The Latent Global Transformer} refers to an autoregressive transformer architecture, $\mathcal{G}$, composed of $l_{\mathcal{G}}$ layers. This model receives a sequence of latent input patch embeddings, denoted as $p_j$, and generates a corresponding sequence of output patch embeddings, $o_j$. In the notation adopted throughout this work, the index $j$ is associated with patches, whereas $i$ indexes bytes. The global model incorporates a block-causal attention mask~\citep{dubey2024llama}, which confines the attention span to patches up to and including the current one within any given document. Notably, this model accounts for the majority of computational operations during both pre-training and inference phases. As such, determining when to execute this model enables dynamic allocation of computational resources based on the complexity of particular input or output segments.

\subsection{Local Encoder}
\label{section:local}

\textit{The Local Encoder Model}, symbolized by $\mathcal{E}$, refers to a compact transformer-based architecture composed of $l_{\mathcal{E}}$ layers, with $l_{\mathcal{E}}$ being significantly fewer than $l_{\mathcal{G}}$. Its central purpose is to rapidly convert an input sequence of bytes $b_i$ into high-dimensional patch embeddings, $p_j$. This model diverges from traditional transformer setups by inserting a cross-attention mechanism after every transformer layer; this addition enables the model to aggregate byte-level features into patch-level embeddings~(\autoref{fig:crossattn}).

Initially, the input bytes $b_i$ are projected into a vector space via a lookup in a $\mathbb{R}^{256 \times h_{\mathcal{E}}}$ embedding table, yielding $x_i$. These vectors can optionally incorporate supplemental context using hash-embeddings~(\S\ref{sec:hash-emb}).

These representations are successively processed through a combination of transformer and cross-attention blocks~(\S\ref{sec:enc-cross-attn}), resulting in the patch representations $p_i$. These patch vectors then serve as input to the global transformer $\mathcal{G}$.

Within each transformer block, attention is restricted using a \textit{local block causal} masking strategy: every byte has access to $w_{\mathcal{E}}$ preceding bytes. This attention window is permitted to span dynamic patch limits, yet always respects document boundaries. In the sections that follow, we detail the embedding construction and the design of the cross-attention modules.

\subsubsection{Encoder Hash n-gram Embeddings}
\label{sec:ngram-emb}
\label{sec:hash-emb}
To build strong and nuanced representations at every position $i$, it is essential to encode contextual information derived from the preceding sequence of bytes. In {\textsc BLT}, this is implemented by representing both the individual byte $b_i$ as well as the sequence-based n-gram containing it. Specifically, for each position $i$, we begin by forming byte-grams

\begin{align}
    g_{i,n}=\{b_{i-n + 1},\ldots, b_i\}
\end{align}

for each byte index $i$ and for values of $n$ ranging from three to eight.\footnote{
Byte-grams of size $n$ or greater are excluded for positions where $i<n$. }

Next, we propose utilizing hash $n$-gram embeddings. In this approach, every byte $n$-gram is assigned, via a hashing function, to an index within a fixed-size embedding table $E_{n}^{hash}$ corresponding to each $n \in \{3,4,5,6,7,8\}$~\citep{bai2010learning}.
These hashed embeddings are then added to the original byte embedding, and the combined representation is subsequently normalized and provided as input to the local encoder. The computation of the enriched embedding is as follows:

\begin{align}
e_i &= x_i + \sum_{n = 3,...,8}E_{n}^{hash}(\text{Hash}(g_{i,n})) \\
\text{where, Hash}(g_{i,n}) &= \text{RollPolyHash}(g_{i,n}) \% |E_{n}^{hash}|
\end{align}

We adjust $e_i$ by dividing it by the number of $n$-gram sizes incremented by one, and implement RollPolyHash as specified in Appendix~\ref{appendix:rollpolyhash}. Section~\ref{section:ablations} explores how varying values of $n$ and different embedding table sizes in $n$-gram hash embeddings influence trends in the flop-controlled scaling law, through ablation studies. Beyond hash-based $n$-gram embeddings, we investigated frequency-based $n$-gram embeddings; further information on these experiments is included in Appendix~\ref{appendix:ngrams}.

\subsubsection{Encoder Multi-Headed Cross-Attention}
\label{sec:enc-cross-attn}
Our approach largely mirrors the input cross-attention mechanism from the Perceiver framework~\citep{jaegle2021perceiver}, with a key distinction: in our model, the latent variables are aligned with patch representations that vary across inputs, rather than employing a uniform set of latent variables~(\autoref{fig:crossattn}). Furthermore, each latent variable exclusively attends to the bytes within its specific patch. Within this module, a query vector is created for every patch $p_j$ by first aggregating its byte-level representations (through pooling) and then applying a linear transformation, $\mathcal{E}_{C} \in \mathbb{R}^{h_{\mathcal{E}} \times (h_{\mathcal{E}} \times U_{\mathcal{E}})}$, where $U_{\mathcal{E}}$ stands for the number of cross-attention heads in the encoder. More specifically, given $f_{\text{bytes}}(p_j)$, which represents the byte sequence belonging to patch $p_j$, we compute

\begin{align}
P_{0,j} &= \mathcal{E}_{C}(f_\text{bytes}((p_j)), f~ \text{is a pooling function} \\
P_l &= P_{l-1} + W_o\left(\text{softmax}\left(\frac{QK^T}{\sqrt{d_k}}\right)V\right) \\
\text{where } Q_j &= W_q(P_{l-1,j}), K_i = W_k(h_{l-1,i}), V_i = W_v(h_{l-1,i}) \\
h_l &= \text{Encoder-Transformer-Layer}_l(h_{l-1})
\end{align}

Here, $P \in \mathbb{R}^{n_p \times h_{\mathcal{G}}}$ denotes a set of $n_p$ patch-level representations that are input to the global processing component. These patch representations are initially formed by aggregating the byte embeddings $e_i$ that are associated with each patch $p_j$. The matrices $W_q$, $W_k$, $W_v$, and $W_o$ correspond to the linear transformations applied to generate the queries, keys, values, and output, respectively; the keys and values are obtained by projecting the byte-level hidden states $h_i$ from the previous network layer (or $e_i$ if it is the initial layer). For attention computation, a specialized patch-based masking approach is applied: each query $Q_j$ is restricted to attend only to the keys and values arising from the bytes within its own patch $j$. Since the model employs multi-headed attention over the $Q$, $K$, and $V$ representations, and considering that the patch embeddings generally possess a higher dimensionality ($h_{\mathcal{G}}$) compared to $h_{\mathcal{E}}$, the patch representations $P_l$ are maintained as separate heads each of dimension $h_{\mathcal{E}}$ during the cross-attention operation; these are subsequently concatenated to recover the full $h_{\mathcal{G}}$ dimensionality. A pre-LayerNorm transformation is applied to the queries, keys, and values, and no positional encodings are introduced in this cross-attention stage. A residual connection also encapsulates the cross-attention module to support gradient flow.

\subsection{Local Decoder} 

Analogous to the local encoder, the local decoder $\mathcal{D}$ is also designed as a compact transformer-based architecture, employing significantly fewer layers ($l_{\mathcal{D}} << l_{\mathcal{G}}$). Its function is to reconstruct sequences of raw bytes, $y_i$, from the global patch representations $o_j$. During this process, the local decoder generates each byte sequentially, conditioning on the sequence of bytes already produced, and leverages the hidden representations generated by the local encoder for the input byte sequence. The decoder alternates through $l_{\mathcal{D}}$ layers that combine cross-attention with standard transformer blocks. Initially, a cross-attention mechanism is utilized to map patch-level representations onto byte-level representations; subsequently, the transformer layer operates on these byte-wise features.

\subsubsection{Decoder Multi-headed Cross-Attention} 

In the decoder cross-attention, the roles of the queries and key/values are interchanged i.e. the byte-representations are now the queries, and the patch representations are now the key/values. The initial byte-representations for the cross-attention are initialized as the byte embeddings from the last encoder layer i.e. $h_{l_{\mathcal{E}}}$. The subsequent byte-representations for layer $l$, $d_{l,i}$ are computed as:

\begin{align}
D_0 &= h_{l_{\mathcal{E}}} \\
B_l &= D_{l-1} + W_o\left(\text{softmax}\left(\frac{QK^T}{\sqrt{d_k}}\right)V\right), \\
  &\text{where } Q_i = W_q(d_{l-1,i}), K_i = W_k(\mathcal{D}_C(o_j)), V_i = W_v(\mathcal{D}_C(o_j)) \\
  D_l &= \text{Decoder-Transformer-layer}_l(B_l)
\end{align}

Here, $W_k$ and $W_v$ serve as key and value projection matrices, respectively, which perform linear transformations and the split operation $\mathcal{D}_C$ on the ultimate patch outputs $o_j$ from the global model. The matrix $W_q$ is utilized as a query projection, acting on the byte-level representations $d_{l-1}$ from the prior decoder transformer layer—or $h_{l_{\mathcal{E}}}$ in the case of the initial layer—while $W_o$ operates as the output projection matrix. This results in $B \in \mathbb{R}^{h_{\mathcal{D}} \times n_b}$, where $n_b$ represents the count of output bytes. Subsequently, the following set of decoder representations $D_l$ is obtained by applying a decoder transformer layer to the cross-attention output $B$. Consistent with the approach taken in local encoder cross-attention, the architecture incorporates multi-head attention, employs pre-LayerNorm, omits positional embeddings, and applies a residual connection around the cross-attention mechanism.

\section{Experimental Setup}
\label{section:experiments}
We established a rigorous experimental framework to systematically assess {\textsc BLT} alongside models utilizing tokenization, ensuring that {\textsc BLT} does not benefit from potential extended sequence contexts that could unfairly affect the comparison.

\subsection{Pre-training Datasets} In all model sizes explored in this work, pre-training is conducted using two datasets: 1) The Llama 2 dataset~\citep{touvron2023llama}, which contains 2 trillion tokens sourced from a range of publicly accessible datasets, followed by thorough data cleaning and filtering steps to enhance data quality; and 2) {\textsc BLT}-1T: A newly constructed corpus consisting of 1 trillion tokens derived from a variety of public data sources, which also incorporates a selection from the Datacomp-LM pre-training dataset~\citep{li2024datacomp}. The Llama 2 dataset serves as the basis for scaling law studies concerning the optimal token count, as recommended by~\cite{dubey2024llama}, which aids in identifying the most effective architectural parameters for {\textsc BLT}. Meanwhile, {\textsc BLT}-1T is employed in a full pre-training process to facilitate direct evaluation against Llama 3 on subsequent tasks. Importantly, none of these datasets involve data collected from Meta's products or services. Additionally, for baseline tokenizer-based model evaluations, the Llama 3 tokenizer with a vocabulary of 128K tokens is utilized, as it has demonstrated superior baseline performance in comparison to the Llama 2 tokenizer in our testing.

\subsection{Entropy Model} The entropy model employed in our experiments is implemented as a byte-level language model, trained on the same data distribution as the complete {\textsc BLT} system. Unless otherwise specified, the architecture used consists of a 100M parameter transformer, comprising 14 layers, with a hidden size of 512 and sliding window attention spanning 512 bytes. All other hyperparameters are aligned with those used in our local and global transformer models. We have conducted experiments exploring various configurations, including different model capacities, receptive field sizes, and model architectures, as detailed in \autoref{section:ablations}. Notably, when the receptive field is sufficiently limited, the trained entropy model can be efficiently represented through a lookup table.

\subsection{Entropy Threshold and Equalizing Context Length} In models that employ entropy-driven patching, we determine a patching threshold aimed at obtaining a target average \textit{patch size} across the pretraining data mixture. Within {\textsc BLT}, in contrast to standard tokenization, the \textit{patch size} can be freely specified, which substantially affects the effective context window available to the model. To prevent models with larger patch sizes from benefiting from longer contexts, we standardize the expected byte count per batch across configurations. Consequently, when patch sizes increase, we proportionally decrease the sequence length to keep the total bytes per batch constant. For Llama 2 data, the context length is fixed at 8k bytes, whereas in the case of the {\textsc BLT}-1T dataset, the context length is expanded to 16k bytes on average, with the batch size consistently set to an average of 16M bytes.

Although the mean batch size remains unchanged, utilizing dynamic patching approaches results in varying proportions of bytes to patches when assembling data batches. To maximize efficiency, our {\textsc BLT} training procedure organizes patches into batches, thereby eliminating the need for padding operations in the resource-intensive latent transformer; this guarantees a uniform count of patches across all batches. During training, we either pad or truncate the byte sequences—setting them to 12k bytes for Llama 2 and 24k bytes for the {\textsc BLT}-1T datasets—to prevent memory usage from spiking due to sequences containing exceptionally large patches.

\subsection{Entropy Model Context}
\label{section:entropy-context}
Our empirical analysis reveals that applying entropy patching results in increasingly extensive patches within structured domains, such as multiple choice questions, which are typically characterized by high repetition (refer to patching on an MMLU instance in \autoref{fig:mmlu_patching}). These patch size increases can be attributed to reduced entropy present in the recurrent sections of the entropy model context. Consequently, for the large-scale application of {\textsc BLT}-Entropy utilizing a patch size of 4.5, we refresh the entropy context by inserting new lines and employ an approximate monotonicity constraint, since this method is less prone to "entropy drift" arising from fluctuations in context length. While this modification alters the manner in which we measure entropies, our strategy for determining the appropriate entropy threshold remains consistent.

\subsection{FLOPs Estimation} 
\label{section:flops}

Our methodology for calculating transformer flops primarily adheres to the framework established in Chinchilla~\citep{hoffmann2022training}, which accounts for the computational costs arising from feed-forward networks, qkvo projections within the self-attention mechanism, and the processes involved in generating attention and output projections. However, in contrast to the original formulation, we presume the input embedding layer utilizes an optimized lookup table rather than a dense matrix multiplication, effectively rendering it a 0-flop computation. Consistent with prior literature, our estimates consider that the backward pass incurs twice the number of flops compared to the forward pass.

To compute flops \textit{per byte} for {\textsc BLT} models, we add up the flops for the local encoder transformer, the global latent transformer, and the local decoder transformer, together with the cross attention blocks in the encoder and the decoder:

\begin{align}
    \text{FL}_{\text{{\textsc BLT}}} &= \text{Transf. FL}(h_{\mathcal{G}}, l_{\mathcal{G}}, m=n_{ctx}/n_p,V=0) / n_p\\
   &+ \text{Transf. FL}(h_{\mathcal{E}}, l_{\mathcal{E}}, m=w_{\mathcal{E}}, V=0) \\
   &+ \text{Transf. FL}(h_{\mathcal{D}}, l_{\mathcal{D}}, m=w_{\mathcal{D}}, V=256) \\ 
    &+ \text{Cross Attn. FL}(h_{\mathcal{E}}, l_{\mathcal{E}}, m=n_p, r=n_p/k) \cdot k/n_p \\ 
   &+ \text{Cross Attn. FL}(h_{\mathcal{D}}, l_{\mathcal{D}}, m=k, r=k/n_p) 
\end{align}

Here, $n_{ctx}$ denotes the sequence length measured in bytes, while $n_p$ represents the patch size. The variable $r$ refers to the proportion of queries relative to keys and values. The parameter $k$ indicates the ratio between the patch dimension and the byte dimension; specifically, it signifies the number of local model partitions that, when concatenated, constitute the full global model representation (for example, $k=2$ in \autoref{fig:crossattn}). The parameter $V$ designates the size of the vocabulary for the output projection layer, utilized exclusively in the local decoder. The module’s operation—whether on the byte sequence or the patch sequence—determines the relevant attention context length, denoted by $m$. We accordingly adjust the calculations for attention FLOPs to reflect each module’s specific configuration. For detailed mathematical formulations for computing FLOPs in Transformer and Cross-Attention modules, refer to Appendix~\ref{section:flops-appendix}.

\subsection{Bits-Per-Byte Estimation} Perplexity only makes sense in the context of a fixed tokenizer as it is a measure of the uncertainty for each token. When comparing byte and token-level models, following previous work~\citep{xue2022byt5, yu2023megabyte, wang2024mambabyte}, we instead report Bits-Per-Byte (BPB), a tokenizer independent version of perplexity. Specifically:

\begin{align}
    \text{BPB}(x)&= \frac{\mathcal{L}_{CE}(\pmb{x})}{\ln(2)\cdot n_{\text{bytes}}}
\end{align}

Here, the aggregate cross-entropy loss, which quantifies the uncertainty inherent in the data $\pmb{x}$, is scaled by dividing by both the total byte count of $\pmb{x}$ and a logarithmic constant.

\subsection{Transformer Architecture Hyperparameters} In constructing all transformer blocks within {\textsc BLT}, encompassing both local and global modules, we largely adopt the design of Llama~3~\citep{dubey2024llama}. Our feed-forward networks utilize the SwiGLU activation~\citep{swiglu}, and rotary positional embeddings (RoPE)~\citep{RoPe} are incorporated exclusively in self-attention layers, employing a rotation parameter of $\theta=500000$~\citep{xiong2024effective}. For normalization, RMSNorm~\citep{RMSNorm} is applied throughout. For self-attention computations with fixed-standard attention masks—such as \textit{block causal} or \textit{fixed-window block causal}—we employ Flash attention~\citep{dao2022flashattention}, configuring a window size of 512 for fixed-width attention types. Because cross-attention in our system utilizes dynamic, patch-dependent masks, we leverage Flex Attention\footnote{\url{https://pytorch.org/blog/flexattention}} to obtain optimized fused kernel implementations, thereby considerably accelerating the training process.

\subsection{{\textsc BLT}-Specific Hyperparameters}  
To evaluate the performance of {\textsc BLT} models, we perform experiments in two main areas: investigating scaling behaviors and assessing performance on downstream tasks. In this context, we use models with various parameter counts: 400M, 1B, 2B, 4B, and 8B. Detailed architecture hyperparameters for these versions can be found in Appendix~Table~\ref{tab:arch}. 

For the first cross-attention layer in the local encoder, we initialize queries via max-pooling. Across all {\textsc BLT} model sizes, the configuration includes $500,000$ hashes utilizing a single hash function and n-gram sizes set between 3 and 8. A constant learning rate of $4e-4$ is employed for every model scale considered. This consistent rate was selected after a hyperparameter sweep within the range of $1e-3$ to $1e-4$ at 400M and 1B, with results indicating that the same rate yielded optimal outcomes for both the token and {\textsc BLT} models. For analyzing scaling on Llama-2 datasets, we follow the training batch-size recommendations from~\cite{dubey2024llama} or apply the equivalent byte-based sizing.

For model optimization, AdamW~\citep{adamw} is used with $\beta_1 = 0.9$, $\beta_2 = 0.95$, and $\epsilon=10^{-8}$. We implement a learning rate schedule consisting of 2000 steps of linear warm-up followed by cosine decay to zero. Additionally, weight decay is applied with a value of 0.1, and global gradient norm clipping is enforced at a threshold of 1.0.

\section{Scaling Trends}
\label{section:scaling}

We provide a comprehensive analysis of scaling behaviors for byte-level models, with implications for the continued development of {\textsc BLT} architectures. This scaling investigation seeks to overcome gaps in earlier studies of byte-level models through the following approaches: (a) we examine the scaling patterns under compute-optimal training conditions, (b) we develop comparable 8B-parameter models using substantial training datasets (reaching up to 1T tokens or 4T bytes) and assess their performance on a range of downstream benchmarks, and (c) we analyze how scaling dynamics manifest under settings where inference costs are fixed. Subsequent sections will delve into particular benefits associated with modeling at the byte-sequence level.

\subsection{Parameter Matched Compute Optimal Scaling Trends}
\label{section:parameter-matched-scaling}
Employing the Llama 2 dataset, we carry out training of both \textit{compute-optimal} bpe and {\textsc BLT} models at four distinct scales, from 1B up to 8B parameters. We then visualize the relationship between the number of training flops and language modeling accuracy on a representative segment of the overall training set. For bpe models, the data-to-parameter ratio is configured according to the guidelines established by Llama~3~\citep{dubey2024llama}, ensuring that the training adheres to the most advantageous model size relative to dataset scale. This \textit{compute-optimal} configuration is motivated by theoretical work indicating it yields maximal performance for the allocated training budget~\citep{hoffmann2022training}, setting a strong comparison point for our approach. Corresponding {\textsc BLT} models are constructed for every bpe model, trained on the identical data subset and employing a Latent Transformer mirroring the architecture and scale of the paired bpe Transformer.

As shown in~\autoref{fig:scaling-compute-opt} (right), {\textsc BLT} models consistently achieve performance on par with or exceeding that of bpe-based models, a pattern that persists across increasing model sizes and greater computational resources. To our knowledge, {\textsc BLT} represents the first byte-level Transformer model to demonstrate scaling behavior comparable to BPE-based architectures at compute-optimal points. These findings support our hypothesis that the optimal relationship between parameter count and training compute established for bpe models similarly extends to {\textsc BLT}, or at least does not diverge significantly.

Both enhancements to architecture and the use of dynamic patching are essential for aligning with bpe scaling behaviors. In ~\autoref{fig:scaling-compute-opt} (left), we present a comparison between models utilizing space-patching and Llama 3. Our estimation of SpaceByte \citep{slagle2024spacebyte} involves implementing {\textsc BLT} space-patching without integrating n-gram embeddings or cross-attention mechanisms. While SpaceByte demonstrates gains over Megabyte, it still does not reach the performance of Llama 3. On the right side of \autoref{fig:scaling-compute-opt}, we depict the cumulative effects arising from architectural modifications combined with dynamic patching. At large scale, {\textsc BLT} models are competitive with leading tokenizer-based models, such as Llama 3.

Additionally, we note that for tokenizer-based models, the selection of tokenizer significantly influences performance; specifically, employing the Llama-3 tokenizer leads to superior results compared to models utilizing the Llama-2 tokenizer when both are trained on identical datasets.

Ultimately, our {\textsc BLT} architecture occupies a middle ground between Llama 2 and Llama 3 when operating with considerably increased patch sizes. The average token sizes produced by the bpe tokenizers for Llama 2 and Llama 3 are 3.7 and 4.4 bytes, respectively. By contrast, {\textsc BLT} demonstrates comparable scaling properties while utilizing average patch sizes of 6 or even 8 bytes. Since inference flop is inversely related to the average patch size, selecting an 8-byte patch size can result in approximately a 50\% reduction in inference computational requirements. Moreover, models with greater patch sizes tend to yield improved outcomes as both the model and dataset are scaled up. For example, while {\textsc BLT} with an 8-byte patch size initially underperforms relative to bpe Llama 2 at the 1B parameter level, it surpasses the bpe baseline when scaled to 7B parameters. This observation indicates that utilizing such large patch sizes could lead to even better results at higher scales, and that it may be practical to explore even larger patch sizes as model capacity and available training compute continue to increase.

\subsection{Beyond Compute Optimal Task Evaluations}
\label{section:evals}
To further investigate the scaling behavior, we extend training of an 8B {\textsc BLT} model past the compute optimal point on the {\textsc BLT}-1T dataset— a more extensive, high-quality corpus— and evaluate its capabilities on a broad set of widely-used classification and generation benchmarks. The evaluation targets common sense reasoning, knowledge, and program synthesis through the following tasks:
\paragraph{Classification tasks} consist of ARC-Easy (0-shot)~\citep{Eval_ARC}, Arc-Challenge (0-shot)~\citep{Eval_ARC}, HellaSwag (0-shot)~\citep{Eval_hellaswag}, PIQA (0-shot)~\citep{Eval_piqa}, and MMLU (5-shot)~\citep{hendrycks2020measuring}. We use a prompt-scoring approach, which involves computing the probability over possible answer choices, and present mean accuracy as the evaluation metric.
\paragraph{Coding related generation tasks:}  To assess the proficiency of LLMs in generating Python solutions, we provide pass@1 results for MBPP (3-shot)~\citep{Eval_MBPP} and HumanEval (0-shot)~\citep{Eval_HumanEval}.

As shown in \autoref{tab:evals}, we evaluate three models that have been trained on the {\textsc BLT}-1T dataset: a model utilizing the Llama 3 tokenizer with byte pair encoding (BPE),\footnote{The Llama 3 tokenizer, featuring a 128k vocabulary, is selected due to its superior performance over the 32k vocabulary of Llama 2.} and two variations of the {\textsc BLT} architecture. The first of these uses a space-based patching approach ({\textsc BLT}-Space), while the second applies an entropy-driven patching method ({\textsc BLT}-Entropy) that incorporates an approximate monotonicity constraint and reinitializes the entropy model's context at new line characters, as outlined in~\autoref{section:entropy-context}. All models are trained with comparable compute in terms of floating point operations. For the {\textsc BLT}-Entropy variant, an additional inference-time modification involves lowering the entropy threshold from 0.6 to 0.1. This adjustment yields better task accuracy but increases the number of inference steps required.

The {\textsc BLT}-Entropy model achieves superior results compared to the Llama 3 model on 4 of the 7 evaluated tasks, even though both models are trained with an equal byte budget. This enhancement in performance can likely be attributed to two main factors: (1) more efficient utilization of training compute made possible by dynamic patching, and (2) the explicit modeling of information at the byte level rather than at the token level.

Conversely, {\textsc BLT}-Space does not match the performance of the Llama 3 tokenizer on most tasks, though it notably decreases inference flops due to its higher average patch size of 6 bytes. In contrast, both bpe and entropy-patching models maintain an average patch size of around 4.5 bytes when evaluated on the same training dataset mixture. Under equivalent training constraints, the model utilizing the larger patch size processes 30\% more data than the others, potentially moving {\textsc BLT} further from compute-optimal efficiency.

\subsection{Patches Scale Better Than Tokens} In {\textsc BLT} models, it is possible to enlarge both the model capacity and the patch size together, all while preserving a constant compute budget for training and inference and without requiring additional training data. This flexible scaling of patch size distinguishes patch-based approaches from fixed-vocabulary, token-based frameworks, as previously addressed in Section~\ref{section:bpe-patch}. Utilizing larger patches reduces computational demands, allowing for those resources to instead be devoted to expanding the global latent transformer, since it needs to be invoked fewer times.

We perform a fixed inference scaling study to evaluate whether larger models executing fewer steps on larger patches surpass the performance of smaller models using more steps. Using initial models of 400m and 3.6B parameters built with the Llama 2 tokenizer, we select flop-matched models using the Llama 3 tokenizer and {\textsc BLT}-Entropy models averaging patch sizes of 6 and 8 bytes on the training datamix (refer to~\autoref{tab:fixed-inf-params} for detailed model specifications). For models with patch size 8, we employ 3 encoder layers instead of 1. Each model is trained with a range of training flop budgets.

As illustrated in \autoref{fig:fixed_inference_scaling}, {\textsc BLT} models demonstrate superior scaling behavior compared to tokenization-based models across both types of inference flop constraints. While BPE models initially outperform at lower training budgets, their advantage diminishes quickly, with {\textsc BLT} surpassing them soon after leaving the compute-optimal point. In real-world scenarios, allocating additional resources during pretraining can yield a model that performs better under a fixed inference budget. A notable instance of this is observed with the 8B parameter class, exemplified by Llama 3.1, which was pretrained on data volumes roughly 100 times greater than the compute-optimal amount for its size.

The point at which {\textsc BLT} surpasses token-based architectures has moved marginally nearer to the compute-optimal regime with the transition to higher flop capacity models—shifting from three times to roughly 2.5 times the optimal compute allocation. Likewise, the patch size 8 variant in the larger flop class demonstrates a more rapid scaling, outpacing alternative models at an earlier stage. As elaborated in Section~\ref{section:parameter-matched-scaling}, increasing the patch size leads to model performance more closely resembling that of BPE-based systems, especially as model sizes grow. This phenomenon can partially be attributed to the proportionally reduced computational effort allocated to the byte-level Encoder and Decoder modules, whose growth in required flops does not keep pace with that of the Latent Transformer. When expanding the parameter count by a factor of 20—from 400M to 8B parameters—{\textsc BLT}'s localized model parameters approximately double. This scaling pattern is notable, since modifications in patch size impact only the computational cost within the patch Latent Transformer, leaving the byte-level modules largely unaffected. As a consequence, the ratio of {\textsc BLT}-Entropy ps=8 parameters to those in the Llama 2 model only slightly increased, shifting from 1.6x to 1.7x as the overall model size was scaled up.

In conclusion, our investigation into patch-length scaling reveals that the {\textsc BLT} architecture, which processes inputs in patches, exhibits superior scaling behavior when both the patch size and model capacity are enlarged together. This favorable scaling tendency appears to hold and even strengthen as model scale increases.

\section{Byte Modeling Improves Robustness} Additionally, we evaluate the robustness of {\textsc BLT} in contrast with token-based counterparts that do not incorporate direct byte-level information, and we introduce a method for converting pretrained token-based models to operate at the byte level.
\label{sec:robustness}

\subsection{Character-Level Tasks}

An initial impetus for developing byte-level models was their inherent resilience to input noise at the byte granularity, as well as their capacity to recognize the internal structure of tokens—challenges that persist for contemporary tokenizer-based approaches. To assess these qualities, we conduct supplementary evaluations utilizing benchmarks designed to test model robustness under noisy input conditions, alongside evaluations that probe character-level sensitivity across English and multiple languages, covering categories such as digits and phonemes. The outcomes of these experiments are summarized in Table~\ref{tab:char_tasks}.

\paragraph{Noisy Data} We generate perturbed variants of the benchmark classification datasets outlined in Section~\ref{section:evals} to assess and compare the resilience of tokenizer-based approaches against that of {\textsc BLT}. Specifically, we utilize five different character-level perturbation methods to systematically alter the text: (a) \textit{AntSpeak}: All characters are capitalized and separated by spaces across the entire text. (b) \textit{Drop}: 10\% of all characters are randomly omitted. (c) \textit{RandomCase}: Each character is assigned to be either uppercase or lowercase at random, with half in each category. (d) \textit{Repeat}: 20\% of the text’s characters are repeated up to four consecutive times. (e) \textit{UpperCase}: The text is entirely converted to uppercase letters. Each noising technique is individually applied to the prompt, the completion, or to both, yielding separate tasks whose results are averaged for evaluation. Table~\ref{tab:char_tasks} presents performance outcomes on noised HellaSwag~\citep{Eval_hellaswag}, revealing that {\textsc BLT} consistently exhibits greater robustness than tokenizer-based models. On average, {\textsc BLT} achieves an 8-point higher score compared to a baseline model trained on equivalent data, and also outperforms the Llama 3.1 model, which has been trained on a considerably larger dataset.

\paragraph{Phonology - Grapheme-to-Phoneme (G2P)} We evaluate the performance of {\textsc BLT} in converting a given sequence of graphemes (i.e., the written characters of a word) into a corresponding phonemic transcription (representing the pronunciation). Table \ref{tab:char_tasks} summarizes the G2P task outcomes under a 5-shot evaluation, utilizing Phonology Bench~\citep{suvarna2024phonologybench}. Our analysis reveals that {\textsc BLT} surpasses the Llama 3 1T tokenizer-based baseline model in this task.

\paragraph{CUTE}
To evaluate character-level processing, we test {\textsc BLT} using the CUTE benchmark~\citep{edman2024cute}, which is composed of multiple tasks divided into three principal areas: compositional understanding, orthographic similarity, and sequence manipulation capabilities. The CUTE benchmark remains particularly difficult for most models that rely on tokenizers. While such models retain certain knowledge about token spellings, they generally struggle to leverage that information when transforming or manipulating text. As presented in Table~\ref{tab:char_tasks}, {\textsc BLT}-Entropy exceeds the performance of both BPE Llama 3 models by over 25 points on this suite of tasks. Notably, our model achieves nearly perfect results—99.9\%—on both tasks that require spelling manipulation, highlighting its effectiveness with character-level changes. This substantial improvement, realized even though {\textsc BLT} was trained on only one-sixteenth the data used for Llama 3.1, implies that BPE-based architectures encounter notable difficulties acquiring character-level competencies. Figure \ref{fig:cute_examples} presents several case studies illustrating contexts in which the Llama 3 tokenizer model falters while {\textsc BLT} maintains robust performance. The only instances where BPE methods surpass our approach are in word insertion and deletion tasks. Since byte-level models do not inherently structure input at the word level, such operations might present additional challenges, though the performance gap remains relatively narrow. Moreover, bottom-up modeling from characters to words may ultimately prove more flexible than attempting the reverse. All experiments adopt a uniform evaluation protocol with prompts sourced directly from Huggingface. Additional prompt optimization could further benefit BPE-based models.

\paragraph{Low Resource Machine Translation} We assess the performance of {\textsc BLT} on translation tasks both to and from six prominent language families, encompassing twenty-one lower resource languages with diverse scripts, utilizing the FLORES-101 benchmark~\citep{goyal2022flores}. SentencePiece BLEU scores are presented in Table~\ref{tab:flores}. The findings indicate that {\textsc BLT} delivers stronger performance than a model built with the Llama 3 tokenizer, recording a 2-point overall improvement when translating into English and a 0.5-point increase when translating out of English. In widely studied language pairs, {\textsc BLT} matches or slightly surpasses Llama 3. Notably, {\textsc BLT} consistently outperforms Llama 3 across multiple low-resource language pairs, highlighting the superiority of byte-based approaches in effectively handling rare and diverse byte sequences.

\subsection{Training {\textsc BLT} from Llama 3}
\label{sec:distillation}
We investigate a method by which {\textsc BLT} models can capitalize on previously trained tokenizer-based architectures to achieve more rapid and efficient convergence during training. Specifically, the approach consists of initializing the global transformer parameters of {\textsc BLT} with those obtained from a pre-trained Llama 3.1 model. Following this initialization, the global transformer’s weights are updated using a learning rate set to one-tenth of that applied to the local encoder and decoder, employing the optimal number of update steps as established for Llama 3. A comparative analysis is provided in Table~\ref{tab:distill} against a baseline {\textsc BLT} model. The results clearly demonstrate that the {\textsc BLT} initialized from Llama 3.1 exhibits substantial gains over both the baseline {\textsc BLT} and the original Llama 3, when normalized by computational resources (measured by number of flops). Notably, even relative to our {\textsc BLT}-Entropy variant (see Table~\ref{tab:evals}), which had been trained on a much larger set of 1T tokens, the {\textsc BLT} initialized from Llama 3.1 achieves superior MMLU performance. This indicates that such initialization constitutes an efficient strategy for markedly reducing training flops while maintaining or improving downstream performance.

This approach effectively reconfigures tokenizer-based architectures into tokenizer-free models, essentially transforming a pre-trained LLaMA 3.1 into a {\textsc BLT} system. For a thorough evaluation, we report results for both the original LLaMA 3.1 model, trained on 15T tokens, and the {\textsc BLT} version derived from it in Table \ref{tab:distill}. While our model shows a minor decrease in performance on MMLU and HumanEval, the reductions are more pronounced on additional benchmarks. These findings indicate that optimizing the adaptation process—such as by refining data mixtures and adjusting hyperparameters—is necessary to fully exploit the capabilities of the pre-trained model and close the performance gap.

\section{Ablations and Discussion}
\label{section:ablations}
Here, we examine various ablation studies that provide rationale for specific architectural decisions in {\textsc BLT}, as well as the selection of the patching method and associated hyper-parameters utilized in training the {\textsc BLT} 8B parameter model on the {\textsc BLT}-1T dataset.

\paragraph{Entropy Model Hyper-parameters} In order to investigate how the size of the entropy model and the length of its context window influence scaling behavior, we conduct experiments with byte-level entropy transformer models ranging from 1m up to 100m parameters, and context window lengths spanning 64 to 512 tokens. For our analysis, we generate bpb versus training flop scaling law plots, using our $400m$ and $1b$ {\textsc BLT} models, both of which are trained on the Llama-2 dataset. The resulting curves are illustrated in \autoref{fig:entropy_ablation}. The results indicate that scaling up either the entropy model’s parameter count or context window tends to enhance performance; however, gains become marginal once the parameter count surpasses 50m.

\paragraph{Types of Patching}

We conduct an ablation study on the four patching approaches described in Section~\ref{section:patching}: 1) Strided Patching with strides of 4 and 6, 2) Whitespace-based Patching, 3) BPE Tokenizer patching utilizing the Llama 3 tokenizer, and 4) Entropy-driven patching employing a compact byte-level LLM.

Although dynamic patching shortens the apparent sequence length, we ensure that context length remains consistent across all patching approaches by adjusting for sequence length. Consequently, each model processes an equivalent expected number of bytes per sequence throughout training and evaluation phases, avoiding potential confounds related to access to longer contexts. The outcomes of these comparative analyses are illustrated in Figure~\ref{fig:scaling-compute-opt}. All alternative patching strategies demonstrate superior performance over static patching, with space patching showing results nearly matching those of dynamic entropy-based patching.

In \autoref{tab:evals_ablation}, we report the results of benchmark evaluations for {\textsc BLT} models, comparing those based on tokenizers, those utilizing space patching, and those employing entropy-based patching. All models were trained for an optimal number of steps on the Llama 2 dataset~\citep{dubey2024llama}. While space patching constitutes a more straightforward method that does not require execution of an entropy model during training, our findings indicate that the improvements attributed to entropy-based patching in scaling analyses~(Section~\ref{section:scaling}) are also evident in performance on downstream benchmark tasks.\footnote{The results for space patching derive from earlier experiments conducted without cross-attention; however, comparable patterns emerge when cross-attention is included.}

\paragraph{Cross-Attention} 
\autoref{tab:crossattn_ablations_1b} presents results from ablation studies examining the impact of incorporating cross-attention at different positions within the encoder and decoder of {\textsc BLT}. For the encoder's cross-attention mechanism, we investigate three distinct query initialization strategies: (1) employing a single learned embedding shared across all global states, (2) generating hash embeddings derived from the patch bytes, and (3) utilizing pooled representations of the encoder's hidden states corresponding to the patch bytes at the specified encoder layer.

Our results indicate that implementing cross-attention within the \textit{decoder} yields the greatest effectiveness. In contrast, incorporating cross-attention into the encoder provides only marginal benefits, and these are observed solely when queries are initialized through pooling. Moreover, the utility of cross-attention becomes more evident on the Common-Crawl dataset, with the effect being most pronounced when employing larger patch sizes.

\paragraph{n-gram Hash Embeddings} 

We evaluate the effect of varying n-gram hash embedding vocabulary sizes (0, 100K, 200K, and 400K) with results summarized in Table~\ref{tab:ngram_ablations_1b}. Our analysis shows that incorporating hash embeddings yields performance improvements across all evaluated domains, with particularly notable benefits on Wikipedia and Github datasets (showing a 0.04 bpb improvement versus a 0.01 bpb gain at 15k steps for an 8B model). When examining a change from 500K to 300K hashes at the 8B scale, performance varied by only approximately 0.001 bpb after 15k steps. These findings suggest that hash embeddings are essential for enabling {\textsc BLT} to attain competitive performance with tokenizer-based models, yet further increases beyond 300K hashes result in minimal additional benefit. Moreover, the observed advantages from hashes seem to operate synergistically with those from cross-attention, as each method contributes to improvements on distinct datasets.

\paragraph{Local Model Hyperparameters} Table \ref{tab:local_layers} presents an ablation study on different configurations concerning the layer count within the local encoder and decoder. Our results indicate that, in conjunction with hash n-gram embeddings, {\textsc BLT} achieves strong performance when the encoder is exceptionally shallow—using only a single layer—while the decoder utilizes a more complex, multi-layer structure.

\section{Related Work}
\label{section:related_work}

\textbf{Character-Level RNNs:} Character Language Modeling has remained an important area of research within neural networks since their inception~\citep{sutskever2011generating,mikolov2012subword,graves2013generating}, primarily due to their capacity to handle out-of-vocabulary words in a natural manner, eliminating the reliance on back-off strategies. \cite{kim2016character} developed a model leveraging purely character-based inputs via convolutional and highway layers, with outputs passed to LSTM RNNs. Their approach achieved performance comparable to then-current state-of-the-art RNN language models on English, while surpassing them for languages with rich morphology—a notable benefit of character-level LLMs. Additionally, \cite{kenter2018byte} applied byte-level LSTM models to machine comprehension tasks, which surpassed word-level models, particularly for highly inflected languages like Turkish and Russian. Along the same trajectory, \cite{zhang2015character} explored character-based convolutional architectures for classification tasks, showing superior performance over word-level baselines in specific scenarios. In further advancements, \citet{chung2019hierarchical} introduced hierarchical LSTM frameworks equipped with boundary detectors at each level, facilitating automatic discovery of latent text hierarchies and yielding enhanced results on character-based language modeling tasks. For machine translation, ByteNet~\cite{kalchbrenner2016neural} adopted CNN layers operating directly on characters, rather than leveraging attention mechanisms.

\textbf{Character-Level Transformers:} The introduction of attention-based transformer architectures~\citep{vaswani2017attention} alongside subword tokenization approaches~\citep{sennrich-etal-2016-neural} has led to marked advancements in neural language modeling and various benchmark evaluations. Nevertheless, utilizing word or subword units embeds a predefined inductive bias, specifying the abstraction layer at which these models operate. Seeking to bridge transformer methods with early, encouraging results in character-level language modeling, \citet{al2019character} explored training extremely deep transformers and leveraged auxiliary loss functions, resulting in transformer-based models that surpassed previous character-level LSTMs. Despite these improvements, their models maintained a considerable performance gap relative to large word-level language models. Moreover, GPT-2~\citep{radfordgpt2} highlighted that for expansive datasets such as the 1 billion word benchmark, language models trained at the byte level still lagged behind those built on word-level units.

\cite{Choe2019BridgingTG} showed that transformer-based language models operating at the byte level can achieve better performance than their subword counterparts when matched for parameter size; however, these models demand significantly more computational resources and require longer training times. Likewise, \cite{el2020characterbert} introduced CharFormer, a BERT variant utilizing convolutional operations over character embeddings to construct word representations, which led to performance gains in medical-domain tasks but similarly incurred higher computational costs. The CANINE model proposed by \cite{clark2022canine} is an encoder-only architecture comprising 150M parameters that processes raw character sequences directly. CANINE's structure features a deep transformer backbone akin to the global model explored here, complemented by a local transformer and strided convolutional layers to perform downsampling on the input characters. CANINE surpasses the performance of comparable token-based encoder-only models (such as mBERT) across a range of multilingual benchmarks. ByT5~\citep{xue2022byt5} investigated the use of byte-level encoder-decoder architectures that forgo patching mechanisms. Although ByT5 achieved heightened robustness to input noise and was competitive with tokenizer-based systems even when trained on a quarter of the data, its omission of patching necessitated costly attention operations over every byte, resulting in substantial computational overhead. Working at the byte level naturally extends sequence lengths compared to subword representations, making it difficult to scale such models efficiently. Most recently, \cite{wang2024mambabyte} leveraged the Mamba Architecture~\citep{gu2023mamba}, which supports a fixed-size memory state over extended contexts, to train a byte-level Mamba model without incorporating patching. This approach yielded superior results compared to byte-level transformer baselines at the 350M parameter scale—specifically in bits-per-byte metrics across multiple datasets—within a fixed computational budget.

\textbf{Patching-based approaches:} Leveraging patching methods can substantially reduce the elevated computation (measured in flops) typically associated with byte-level LLMs, all while maintaining competitive performance. Several studies have achieved promising results on smaller-scale models and modest training data. In particular, \cite{nawrot-etal-2022-hierarchical} conducted experiments incorporating fixed patching-based downsampling and upsampling within their hourglass transformer, which exhibited superior results compared to alternative byte-level baselines at the 150M parameter scale. Building upon this, \cite{nawrot-etal-2023-efficient} introduced advances by utilizing dynamic patching strategies. These involve a boundary-predictor trained end-to-end, a version guided by tokenizer supervision, and an entropy-driven patching approach akin to {\textsc BLT}. Their findings indicated that these strategies surpass standard transformers in language modeling tasks, particularly with models around 40M parameters on datasets of approximately 400M tokens. Moreover, \cite{lester2024training} explored model training on sequences compressed through arithmetic coding to achieve higher compression rates than those attainable with BPE. By employing an equal-info windows approach, they managed to improve upon byte-level baselines in language modeling scenarios, though subword-based baselines still performed better.

Our research builds upon and is closely aligned with the work on MegaByte~\citep{yu2023megabyte}, a causal large language model (LLM) that relies on a decoder-only framework. MegaByte leverages a fixed static patching approach: byte sequences are divided into patches by concatenating their representations, after which a local model within the decoder reconstructs the patches back into individual bytes. Their results indicate that MegaByte attains performance comparable to tokenizer-based approaches at the 1B parameter scale on data comprising roughly 400B bytes. In our study, we systematically examine MegaByte's performance and observe that static patching methods tend to underperform relative to contemporary compute-optimal tokenizer-based models when controlling for floating-point operations (flops); our results further highlight how {\textsc BLT} addresses this performance discrepancy. Similarly, \cite{slagle2024spacebyte} reach analogous conclusions regarding MegaByte’s limitations. They propose modifications, such as applying patching to whitespaces and other space-resembling bytes, alongside introducing a local encoder. Their experimental results yield improvements over token-based transformers in certain domains (notably Github and arXiv) at the 1B scale under compute constraints. We replicate experiments with this variant and demonstrate that, despite gains, further architectural innovations remain necessary for byte-level models to achieve true parity with current state-of-the-art models like Llama 3.

\section{Limitations and Future Work}

For the selection of architectural parameters in this study, we adopt the number of training steps established as optimal for Llama~3~\citep{dubey2024llama}. Notably, these scaling principles were derived for BPE-level transformers, and applying them to {\textsc BLT} may result in less-than-ideal data-to-parameter ratios. Developing precise scaling laws tailored to {\textsc BLT} remains an open task and could potentially reveal even more advantageous scaling dynamics specific to our architecture. Furthermore, a significant portion of the empirical evaluations were performed at parameter scales up to 1B; it remains an open question whether the most effective architectural configurations will shift as models are extended to 8B parameters or more, potentially unlocking superior performance at these higher scales.

Current transformer frameworks and code repositories are optimized primarily for architectures built around tokenizers. Although we conduct experiments that match theoretical FLOPs and make use of some efficient methods like FlexAttention for components that depart from the standard transformer setup, there may still be discrepancies in wall-clock performance compared to tokenizer-based systems. Consequently, additional optimization could further improve the speed and efficiency of our implementations.

Although {\textsc BLT} currently employs an independently trained entropy model for the patching process, a promising avenue for future research lies in jointly optimizing the patching model in an end-to-end manner. In Section \ref{sec:distillation}, we offer preliminary findings that suggest it is possible to ``byte-ify'' large tokenizer-based models—such as Llama 3, which is trained on over 10T tokens—by initializing the global transformer with their parameters and keeping those parameters fixed. Advancing this line of work could yield techniques that preserve the advantages of bytefying while potentially surpassing the performance of existing tokenizer-based models, all without the need to retrain them from the beginning.

\section{Conclusion}
\label{section:conclusion}

This work introduces the Byte Latent Transformer (\textbf{BLT}), an innovative framework that challenges the traditional reliance on fixed-vocabulary tokenization within large language models. By adopting a dynamic and trainable mechanism for aggregating bytes into variable-size patches, {\textsc BLT} enables adaptive allocation of computational effort according to the intrinsic complexity of the input data. This results in notable gains in both computational efficiency and resilience to challenging inputs. Through a comprehensive scaling analysis, we show that {\textsc BLT} is capable of achieving performance levels comparable to established tokenization-driven models such as Llama 3 at scales of up to 8B parameters and 4T bytes. Importantly, BLT models can achieve as much as a 50\% reduction in inference floating-point operations by accepting slight decreases in standard evaluation scores. In addition, {\textsc BLT} facilitates novel scaling strategies, such as jointly enlarging model parameters and patch sizes without exceeding a fixed inference budget, a crucial property for real-world computational settings. By processing data directly at the byte level, {\textsc BLT} also demonstrates enhanced proficiency in managing rare or noisy input cases and provides improved modeling of sub-word phenomena. In summary, the evidence presented here positions {\textsc BLT} as a strong contender to traditional tokenization approaches, delivering a flexible, efficient, and resilient architecture for the development of large language models.

\end{document}